\documentclass[11pt]{article}

\usepackage[a4paper,margin=1in]{geometry}
\usepackage{fontspec}
\usepackage{xcolor}
\usepackage{hyperref}
\usepackage{enumitem}
\usepackage{amsmath}
\usepackage{booktabs}
\usepackage{array}

\setmainfont{Helvetica Neue}
\setsansfont{Helvetica Neue}
\setmonofont{Menlo}
\newfontfamily\sinhalafont[Script=Sinhala]{Sinhala MN}
\newcommand{\si}[1]{{\sinhalafont #1}}

\definecolor{accent}{HTML}{255A73}
\definecolor{softgray}{HTML}{F5F7F8}
\definecolor{darkgray}{HTML}{333333}

\hypersetup{
  colorlinks=true,
  linkcolor=accent,
  urlcolor=accent
}

\setlist[itemize]{leftmargin=1.2em}
\setlist[enumerate]{leftmargin=1.4em}
\renewcommand{\arraystretch}{1.2}

\title{\textbf{Sinhala Stem/Rest Segmentation}\\
\large Split Logic and Scoring Heuristics}
\author{Sarves}
\date{\today}

\begin{document}
\maketitle

\section*{Purpose}

The Sinhala Stem Segmenter is designed for modern written Sinhala. It is not a
lemmatizer, and it does not delete suffixes. Instead, it segments a word into a
left-side \textit{stem} and a right-side \textit{rest}:

\[
\text{stem} + \text{rest} = \text{normalized word}
\]

For example:

\begin{center}
\large \si{ජනාධිපතිවරයා} + \si{ය} = \si{ජනාධිපතිවරයාය}
\end{center}

The goal is to preserve the original word while exposing a useful stem/rest
structure for search, tokenization experiments, corpus analysis, and downstream
Sinhala NLP tools.

\section*{What Happens When A Sinhala Word Is Fed?}

When a word such as \si{ජනාධිපතිවරයාය} is given to the segmenter, the model
does not simply look up a final answer from a dictionary. It uses stored learned
evidence, but the split itself is generated and scored on the fly.

\begin{enumerate}
  \item \textbf{Normalize the word.}
  The input is Unicode-normalized so equivalent written forms are handled
  consistently.

  \item \textbf{Split into Sinhala orthographic clusters.}
  The word is not split by raw Unicode characters. The segmenter respects the
  Sinhala abugida writing system and avoids cutting inside dependent vowel
  signs, al-lakuna, ZWJ/ZWNJ sequences, and related written units.

  \item \textbf{Generate legal split candidates on the fly.}
  For \si{ජනාධිපතිවරයාය}, the segmenter considers legal boundaries such as:

  \begin{center}
  \large \si{ජනාධිපති} + \si{වරයාය}
  \end{center}

  \begin{center}
  \large \si{ජනාධිපතිවරයා} + \si{ය}
  \end{center}

  and other legal cluster-safe splits.

  \item \textbf{Check learned evidence.}
  The model contains learned tables such as:
  \begin{itemize}
    \item stem candidate $\rightarrow$ possible rest parts and frequencies
    \item rest part $\rightarrow$ global frequency
    \item word $\rightarrow$ frequency
  \end{itemize}

  \item \textbf{Score each candidate split.}
  Each legal split is scored using corpus evidence.

  \item \textbf{Choose the best split.}
  The highest-scoring valid split is selected if it is above the minimum score
  threshold.

  \item \textbf{Fallback if no split is strong enough.}
  If the model does not find a strong split, it returns the whole word as the
  stem and an empty rest.
\end{enumerate}

\section*{Example Output}

For the input word \si{ජනාධිපතිවරයාය}, the model returns:

\begin{center}
\begin{tabular}{>{\bfseries}p{0.28\linewidth}p{0.52\linewidth}}
\toprule
Field & Value \\
\midrule
word & \si{ජනාධිපතිවරයාය} \\
normalized & \si{ජනාධිපතිවරයාය} \\
stem & \si{ජනාධිපතිවරයා} \\
rest & \si{ය} \\
boundary & 12 \\
stem\_clusters & 8 \\
rest\_clusters & 1 \\
confidence & 1.0 \\
score & 42.895175 \\
family\_size & 76 \\
family\_tokens & 22839 \\
rest\_frequency & 7212294 \\
split & True \\
\bottomrule
\end{tabular}
\end{center}

The important user-facing result is:

\begin{center}
\large \si{ජනාධිපතිවරයාය} \(\rightarrow\) \si{ජනාධිපතිවරයා} + \si{ය}
\end{center}

\section*{Meaning Of The Advanced Fields}

\begin{description}
  \item[word] Original input word.
  \item[normalized] Unicode-normalized word used internally.
  \item[stem] Left-side segment selected by the model.
  \item[rest] Right-side remaining segment.
  \item[boundary] Character offset where the stem ends and the rest starts.
  \item[stem\_clusters] Number of Sinhala orthographic clusters in the stem.
  \item[rest\_clusters] Number of Sinhala orthographic clusters in the rest.
  \item[confidence] Model confidence between 0 and 1.
  \item[score] Raw internal evidence score for the chosen split.
  \item[family\_size] Number of different rest parts seen with the same stem.
  \item[family\_tokens] Total frequency supporting this stem family.
  \item[rest\_frequency] Frequency of this rest pattern in the model.
  \item[split] True if a split was selected; False if the whole word was kept as
  the stem.
\end{description}

\section*{Scoring Formula}

For each legal candidate split, the model calculates:

\[
\begin{aligned}
\text{score} ={}&
  2.2 \cdot \log(1 + \text{family\_size}) \\
&+ 0.65 \cdot \log(1 + \text{family\_tokens}) \\
&+ 1.15 \cdot \log(1 + \text{rest\_frequency}) \\
&+ 0.7 \cdot \log(1 + \text{stem\_as\_word\_frequency}) \\
&+ 0.12 \cdot \text{stem\_clusters} \\
&- 0.45 \cdot \max(0, \text{rest\_clusters} - 4) \\
&- \text{tiny\_stem\_penalty}
\end{aligned}
\]

where:

\[
\text{tiny\_stem\_penalty} =
\begin{cases}
0.7, & \text{if stem\_clusters} \leq \text{min\_stem\_clusters} \\
0, & \text{otherwise}
\end{cases}
\]

\section*{Logic Behind The Heuristics}

The formula is a tunable engineering heuristic. It is not copied from a single
published Sinhala stemming paper. However, the signals used in it are motivated
by established ideas in unsupervised morphology and stemming.

\subsection*{Family Size}

\[
2.2 \cdot \log(1 + \text{family\_size})
\]

This is the strongest signal. If a left-side string appears with many different
right-side parts, it is likely to be a stem-like unit.

Example pattern:

\begin{center}
\large \si{ආණ්ඩුව}+\si{ට},\quad
\si{ආණ්ඩුව}+\si{ේ},\quad
\si{ආණ්ඩුව}+\si{ක්}
\end{center}

This idea is related to \textit{successor variety}: a high variety of possible
continuations after a segment can indicate a boundary.

\subsection*{Family Tokens}

\[
0.65 \cdot \log(1 + \text{family\_tokens})
\]

This rewards stem families that occur often. It is weighted lower than
family size because raw frequency alone can over-favor very common accidental
strings.

\subsection*{Rest Frequency}

\[
1.15 \cdot \log(1 + \text{rest\_frequency})
\]

This rewards common right-side patterns. If a rest part appears across many
words, it is a useful segmentation clue.

\subsection*{Stem As Word Frequency}

\[
0.7 \cdot \log(1 + \text{stem\_as\_word\_frequency})
\]

This rewards stems that also appear independently as words. This is only a soft
clue; a useful stem does not have to be a standalone dictionary word.

\subsection*{Stem Length Reward}

\[
0.12 \cdot \text{stem\_clusters}
\]

This gives a small reward to longer stems, helping avoid splits that happen too
early in the word.

\subsection*{Long Rest Penalty}

\[
-0.45 \cdot \max(0, \text{rest\_clusters} - 4)
\]

This penalizes unusually long rest parts. A very long rest can indicate that the
split happened too early.

\subsection*{Tiny Stem Penalty}

Very short stems are penalized because they are more likely to be accidental
prefix-like fragments rather than useful stems.

\section*{Why Use Logarithms?}

The model uses \(\log(1+x)\) for count-based signals. This gives diminishing
returns:

\begin{itemize}
  \item frequency evidence helps;
  \item very large counts do not completely dominate the decision;
  \item multiple kinds of evidence can still influence the split.
\end{itemize}

\section*{Literature Support}

The exact numeric weights are engineering choices and should be tuned with
evaluation data. The underlying ideas, however, have support in computational
morphology:

\begin{itemize}
  \item \textbf{Successor variety:} Harris-style segmentation uses the idea that
  many possible continuations after a segment can indicate a morpheme boundary.
  This supports the use of \textit{family\_size}.

  \item \textbf{Unsupervised morphology:} Morfessor-style models show that useful
  morph-like units can be induced from raw text using frequency and form/usage
  evidence.

  \item \textbf{Conservative stemming:} Stemming for information retrieval often
  prefers avoiding harmful over-stemming. This motivates penalties for tiny stems
  and very long rest parts.
\end{itemize}

\section*{Practical Interpretation}

For most users, the core output is:

\[
\text{word} \rightarrow \text{stem} + \text{rest}
\]

The score and frequency fields are best treated as diagnostic evidence. They
help explain why the model selected a split, but they are not grammatical labels
and not proof that the split is the only linguistically correct analysis.

\section*{References}

\begin{itemize}
  \item Harris, Z. S. (1955). \textit{From phoneme to morpheme}. Language, 31(2),
  190--222.

  \item Çöltekin, Ç. (2010). \textit{Improving Successor Variety for Morphological
  Segmentation}.

  \item Creutz, M., \& Lagus, K. (2007). \textit{Unsupervised models for morpheme
  segmentation and morphology learning}. ACM Transactions on Speech and
  Language Processing, 4(1).

  \item Snowball stemming project. \url{https://snowballstem.org/}
\end{itemize}

\end{document}
